\begin{lemma}
Let \(0<\varepsilon\leq 1\), let \(0\leq p\leq 1\), and let \(m\) be a
natural number such that
\[
  \frac{100}{\varepsilon}\leq p\,m .
\]
Put \(t=\log(1+\varepsilon/2)\).  Then \(t>0\), and the exponential envelope
appearing on the right-hand side of the Bernoulli MGF bound satisfies
\[
\begin{aligned}
m\Big(&\exp\{-t(1+\varepsilon)pm
      +(m-1)p(\exp(t)-1)\}\\
&+\exp\{t(1-\varepsilon)pm
      +(m-1)p(\exp(-t)-1)\}\Big)
\leq
m\Big(\exp\{-(\varepsilon^2/100)pm\}
     +\exp\{-(\varepsilon^2/100)pm\}\Big).
\end{aligned}
\]
\end{lemma}
\begin{proof}
Set \(M=m\), \(q=pM\), \(c=\varepsilon^2/100\), and
\(t=\log(1+\varepsilon/2)\).  Since \(\varepsilon>0\), we have \(t>0\).
Also \(\exp(t)-1=\varepsilon/2\) and
\(\exp(-t)-1=-\varepsilon/(\varepsilon+2)\).

For the upper-tail exponent, use \(m-1\leq m\) and
\(\exp(t)-1\geq0\).  The exponent is at most
\[
\bigl(-(1+\varepsilon)t+\varepsilon/2\bigr)q.
\]
The standard inequality
\[
\frac{2x}{x+2}\leq \log(1+x)\qquad (x\geq0)
\]
with \(x=\varepsilon/2\) gives
\[
(1+\varepsilon)t-\varepsilon/2\geq \varepsilon^2/100
\]
for \(0<\varepsilon\leq1\).  Hence the upper exponent is at most \(-cq\).

For the lower-tail exponent, the assumption \(100/\varepsilon\leq q\) and
\(p\leq1\) imply \(p\leq(\varepsilon/100)q\), so
\[
(m-1)p=q-p\geq(1-\varepsilon/100)q.
\]
Since \(\exp(-t)-1\leq0\), the lower exponent is at most
\[
\left(\frac{\varepsilon}{2}(1-\varepsilon)
  -(1-\varepsilon/100)\frac{\varepsilon}{\varepsilon+2}\right)q.
\]
For \(0<\varepsilon\leq1\), the coefficient in parentheses is at most
\(-\varepsilon^2/100=-c\).  Thus the lower exponent is also at most \(-cq\).
Monotonicity of the exponential function gives the two displayed exponential
bounds, and multiplying their sum by the nonnegative factor \(m\) gives the
claim.
\end{proof}
